\documentclass[12pt,a4paper]{article}

% -------------------------------------------------
% Packages
% -------------------------------------------------
\usepackage[T1]{fontenc}
\usepackage[french]{babel}
\usepackage{lmodern}
\usepackage{amsmath,amssymb,amsfonts}
\usepackage[a4paper,margin=2cm]{geometry}
\usepackage[most]{tcolorbox}
\usepackage{tikz}
\usetikzlibrary{positioning, calc}

% -------------------------------------------------
% Document
% -------------------------------------------------
\begin{document}

\section*{I. Fonction dérivable en un point}

\begin{tcolorbox}[
    colback=green!5,
    colframe=green!60!black,
    colbacktitle=green!60!black,
    title=1. Définition du nombre dérivé,
    fonttitle=\bfseries,
    boxrule=0.8pt
]

Soit $f$ une fonction numérique et $a$ un nombre réel tel que $f$
est définie sur un intervalle $I$ contenant $a$. On dit que $f$
est dérivable en $a$ si :

\[
\lim_{x\to a}\frac{f(x)-f(a)}{x-a}=\ell\in\mathbb{R}
\]

$\ell$ est appelé nombre dérivé de $f$ en $a$ que l'on note
$f'(a)=\ell$.

\end{tcolorbox}

\textbf{Exemples :} Les fonctions suivantes sont-elles dérivables
aux points indiqués ?

\[
f(x)=x^2 \quad \text{en } a=-1
\qquad\qquad
g(x)=2x^3-1 \quad \text{en } a=2
\]

\textbf{Contre exemple :} $f(x)=|x-1|$ et $f(1)=0$.
$f$ est-elle dérivable en $1$ ?

\begin{tcolorbox}[
    colback=green!5,
    colframe=green!60!black,
    colbacktitle=green!60!black,
    title=\textbf{2. Théorème},
    fonttitle=\bfseries,
    boxrule=0.8pt
]

Si une fonction $f$ est dérivable en un réel $a$, alors $f$ est
continue en $a$. C'est-à-dire que la dérivabilité entraîne
la continuité.

\textbf{NB :} La réciproque n'est pas toujours vraie.

\end{tcolorbox}

\textbf{Exemple :}

\begin{enumerate}

    \item Soit
    \[
    g(x)=\frac{f(x)-f(1)}{x-1}
    \]
    avec $f(x)=|x-1|$. $g$ est-elle continue en $1$ ?

    \item Soit $f(x)=|x|$. $f$ est-elle dérivable en $0$ ?

\end{enumerate}

% -------------------------------------------------
% 3. Interprétation graphique
% -------------------------------------------------

\begin{tcolorbox}[
    colback=green!5,
    colframe=green!60!black,
    colbacktitle=green!60!black,
    title=\textbf{3. Interprétation graphique},
    fonttitle=\bfseries,
    boxrule=0.8pt
]

Soit $f$ une fonction numérique dérivable en un réel $a$.
La courbe représentative de la fonction $f$ admet une tangente
d'équation :

\[
(T):\quad y=f'(a)(x-a)+f(a)
\quad \text{au point } A(a;f(a))
\]

Le coefficient directeur de $(T)$ est $f'(a)$.

\end{tcolorbox}

\textbf{Exemples :}

\begin{enumerate}

    \item $f(x)=x^2-3x-1,\qquad x_0=-2$

    \item $f(x)=\dfrac{2-3x}{x-2},\qquad x_0=0$

    \item $f(x)=\sqrt{2x-3},\qquad x_0=2$

\end{enumerate}

% -------------------------------------------------
% 4. Dérivée à gauche et dérivée à droite
% -------------------------------------------------

\begin{tcolorbox}[
    colback=green!5,
    colframe=green!60!black,
    colbacktitle=green!60!black,
    title=\textbf{4. Dérivée à gauche et dérivée à droite},
    fonttitle=\bfseries,
    boxrule=0.8pt
]

Soit $f$ et $a$ un nombre réel :

\begin{itemize}

    \item Si
    \(
    \lim\limits_{x\to a^-}\frac{f(x)-f(a)}{x-a}
    =\ell\in\mathbb{R}
    \)
    alors on dit que $f$ est dérivable à gauche de $a$ et le réel
    $\ell$ est appelé nombre dérivé de $f$ à gauche de $a$ que l'on
    note $f'_g(a)=\ell$.

    \item Si
    \(
    \lim\limits_{x\to a^+}\frac{f(x)-f(a)}{x-a}
    =\ell'\in\mathbb{R}
    \)
    alors on dit que $f$ est dérivable à droite de $a$ et que
    $\ell'$ est le nombre dérivé de $f$ à droite de $a$ que l'on
    note $f'_d(a)=\ell'$.

\end{itemize}

\end{tcolorbox}

\end{document}